\begin{figure}[htbp]
    \centering
    \begin{tikzpicture}[
        node distance=1.6cm,
        % 定义样式以匹配你之前的需求
        simplebox/.style={draw, rectangle, thick, rounded corners=2pt, align=center, inner sep=2pt, fill=white},
        stdarrow/.style={->, >=latex, thick},
        arrowlabel/.style={font=\footnotesize, align=center, midways, above}
    ] 

        % --- Nodes ---
        
        % Node 1: DL Systems
        \node[simplebox, minimum height=1cm] (dlsys) {
            DL Systems
        };

        % Node 2: Context Generation
        \node[simplebox, right=of dlsys, minimum height=1cm] (ctxgen) {
            Context\\[-2pt]Generation\\(\S\ref{sec:design:context_construction})
        };

        % Node 3: Information Propagation
        \node[simplebox, right=of ctxgen, minimum height=1cm] (host) {
            Information\\[-2pt]Propagation\\(\S\ref{sec:design:early-return} \& \S\ref{sec:design:propagation})
        };

        % Node 4: Kernel Specializer
        \node[simplebox, right=of host, minimum height=1cm] (spec) {
            Kernel\\[-2pt]Specializer\\(\S\ref{sec:design:specialization})
        };

        % Node 5: Final Output - 改为纯文本，去掉方框 [Requirement 1]
        \node[right=0.8cm of spec, align=left, font=\small\bfseries] (final) {
            Specialized Kernels\\or Debloated Library
        };


        % --- Connections ---

        % Arrow 1
        \draw[stdarrow] (dlsys) -- node[above, font=\footnotesize, align=center] {Model-specific\\Constants} (ctxgen);

        % Arrow 2
        \draw[stdarrow] (ctxgen) -- node[above, font=\footnotesize, align=center] {Calling\\Context} (host);

        % Arrow 3
        \draw[stdarrow] (host) -- node[above, font=\footnotesize, align=center] {Constant\\Schedule\\Parameters} (spec);

        % Final arrow to output
        \draw[stdarrow] (spec) -- (final);

    \end{tikzpicture}
    % [Requirement 2] 增加 Caption 和 Label
    \caption{The overall workflow of \mysys.}
    \label{fig:overall_workflow}
    \vspace{-10pt}
\end{figure}